% ===== 原文 pp.88–91（PDF p.98–101） =====
% [原文 p.88]
\section{修正 2d A-模型的场含量、作用量与变分}

现在我们把靶流形特殊化为 Kähler 流形 $X$，它带有群 $U(1)$ 的一个左作用；假定该作用保持 $X$ 的辛形式与复结构，即
\[
\mathcal{L}_V\Omega=\mathcal{L}_VJ_X=0,
\]
其中 $V\in\Gamma(TX)$ 是 $U(1)$ 作用的生成向量场。由此可知 $U(1)$ 通过等距作用，因为上式蕴含 $\mathcal{L}_Vg=0$。在这些假定下，A-模型各场上诱导的作用 $U(1)_X$ 保持作用量不变；然而出于稍后会变得明显的原因，我们附加假定 $U(1)$ 的作用是\emph{Hamilton 作用}，即它由 $X$ 上一个整体定义的等变矩映射生成。这句话的意思是：存在整体定义的标量函数 $H:X\to\R$，使得在 $X$ 的每一点都有
\begin{equation}
\dd H=-\iota_V\Omega \tag{7.5}
\end{equation}
其中 $\iota_V$ 表示对向量 $V$ 缩并。按分量写就是 $\partial_JH=\Omega_{IJ}V^I$。由
\[
\dd\iota_V\Omega=\mathcal{L}_V\Omega-\iota_V\dd\Omega=0
\]
可知 1-形式 $\iota_V\Omega$ 自动闭，从而只要 $\Omega$ 闭且在 $V$ 下不变，就存在满足 (7.5) 的局部定义函数；因此 Hamilton 假定中非平凡的部分是函数 $H$ 整体有定义。此外，对群 $U(1)$ 而言矩映射 $H$ 的等变性自动成立，因为
\[
\iota_V\dd H=-\iota_V\iota_V\Omega=0 .
\]

% [Tutor Note: p.88 关于矩映射存在性的论证值得补全：$\iota_V\Omega$ 闭（由 $\dd\iota_V\Omega=\mathcal{L}_V\Omega=0$）只保证它在每个单连通 chart 上恰当（Poincaré 引理），$H$ 整体存在则是上同调条件 $[\iota_V\Omega]=0\in H^1_{dR}(X;\R)$。本段后面所有推导都只用 $H$ 的微分性质，故 $H$ 差一个局部常数是允许的；这个常数自由度在 p.94 又被边界条件（$H$ 在边界的每个分支上取常值）进一步约束。]
% [Tutor Note: p.88 有一处符号需要当心：(7.5) 印作 $\dd H=-\iota_V\Omega$，而同页紧接着的分量式写 $\partial_JH=\Omega_{IJ}V^I$；由 p.86 的约定 $\Omega=\tfrac12\Omega_{IJ}\dd\varphi^I\wedge\dd\varphi^J$ 直接计算得 $\iota_V\Omega=\tfrac12\Omega_{IJ}(V^I\dd\varphi^J-V^J\dd\varphi^I)=\Omega_{IJ}V^I\dd\varphi^J$，故 $\dd H=+\iota_V\Omega$ 才与分量式一致。判断依据在 p.100：该处要由 $S_1=-t\int_\Sigma H(\Phi)\dd\omega$ 的 BRST 变分抵消 $\delta_QS_0=+t\int_\Sigma\Omega_{IJ}V^I\chi^J\dd\omega$，这要求 $\delta_QH=\chi^I\partial_IH=+\Omega_{IJ}V^I\chi^J$，即取正号。译文照录原式的负号，读者自行按此更正即可（p.92 的 $\dd_C\eta=0$ 沿用了负号约定，见该页批注）。]

\para{修正 A-模型的场含量}

如上所述，我们希望按式 (7.1) 修改场的自旋来扭转理论。扭转之后，玻色场 $\Phi$ 不再是简单的从 $\Sigma$ 到 $X$ 的映射，因为它在 $\Sigma$ 的旋转下非平凡地变换。确切地说，它被替换为相伴纤维丛
\[
E=P\times_{U(1)}X
\]
的一个截面 $\Phi$；该丛以 $X$ 为典型纤维，由群作用 $l$ 相伴于 $\Sigma$ 上的主 $U(1)$ 标架丛 $P$。也就是说，$\Phi$ 是从 $\Sigma$ 到全空间 $E$ 的映射，并且是投影映射 $\pi_E:E\to X$ 的一个右逆。同样，扭转理论中的费米场取值于拉回向量丛 $\Phi^*VE$，其中 $VE=\ker\dd\pi_E\subset TE$ 是沿纤维指向的垂直向量构成的子丛。正如附录 C 所讨论的，该丛从 $X$ 的结构自然继承一个

% [原文 p.89]
复结构 $J^V:VE\to VE$，因此把垂直向量投影到其 $(1,0)$ 与 $(0,1)$ 部分是合理的。

总之，修正 A-模型的玻色场由一个截面
\[
\Phi\in\Omega^0(\Sigma,E)
\]
以及辅助场
\begin{gather*}
P\in\Omega^{0,1}\big(\Sigma,\Phi^*V^{1,0}E\big)\\
\tilde P\in\Omega^{1,0}\big(\Sigma,\Phi^*V^{0,1}E\big)
\end{gather*}
组成。修正 A-模型的费米场由形式
\begin{gather*}
\chi\in\Omega^0\big(\Sigma,\Phi^*VE\big)\\
\rho\in\Omega^{0,1}\big(\Sigma,\Phi^*V^{1,0}E\big)\\
\tilde\rho\in\Omega^{1,0}\big(\Sigma,\Phi^*V^{0,1}E\big)
\end{gather*}
组成。关于 $E$ 及其附带结构的几何，详见附录 C。

与第 3 章讨论的带规范场 4d A-模型一样，截面 $\Phi:\Sigma\to E$ 局部地由映射
\[
\Phi^{(\alpha)}:U^{(\alpha)}\to X
\]
表示，而截面 $\chi\in\Omega^0(\Sigma,\Phi^*VE)$ 局部地由截面
\[
\chi^{(\alpha)}\in\Omega^0\big(U^{(\alpha)},\Phi^{(\alpha)*}TX\big)
\]
表示，其中 $U^{(\alpha)}\subset\Sigma$ 是底层主丛 $P$ 的平凡化邻域。因此，局部地看 $\Sigma$ 上修正 A-模型的场含量与普通 A-模型等价。

% [Tutor Note: p.88 把 $\Phi$ 从“映射 $\Sigma\to X$”改成“$E=P\times_{U(1)}X$ 的截面”，这是全章的技术核心，可用局部平凡化核对：在 $U^{(\alpha)}$ 上选 $P$ 的局部截面 $s^{(\alpha)}$，则 $E|_{U^{(\alpha)}}\cong U^{(\alpha)}\times X$，截面对应映射 $\Phi^{(\alpha)}$；在重叠 $U^{(\alpha)}\cap U^{(\beta)}$ 上两个映射相差一个 $U(1)$ 作用 $\Phi^{(\beta)}=l_{(\beta\alpha)}\cdot\Phi^{(\alpha)}$。因为 $U(1)$ 保持 $\Omega$ 与复结构，局部写下的 A-模型公式在换 chart 后只差一个 $U(1)_X$ 变换，这正是 p.85 所说“差别只在 chart 重叠处”。注意 $E$ 的总空间实维数为 $\dim_\R X+1$，但截面仍有 $\dim_\R X$ 个自由度。]

\para{修正 A-模型的作用量}

(7.2) 中离壳形式 BRST 变分的局部表达式在扭转后仍然成立。然而作用量的表达式需要修改，以确保在扭转理论中它在 $\Sigma$ 上整体良定义并且 BRST 不变。为使作用量在 $\Sigma$ 上整体良定义，我们必须把普通导数换成作用在 $E$ 与 $\Phi^*VE$ 的截面上的协变导数。这又要求我们在底层主丛 $P$ 上取定一个联络。我们取 $P$ 上唯一与度规相容且无挠的自旋联络 $\omega$；事实上，我们把这一选择视为“修正 A-模型”定义的一部分。在我们所用的 $\Sigma$ 适配

% [原文 p.90]
局部坐标下，自旋联络显式地由下式给出
\begin{equation}
\omega=(\partial_2f)\dd\sigma^1-(\partial_1f)\dd\sigma^2=\ii(\partial_zf)\dd z-\ii(\partial_{\bar z}f)\dd\bar z . \tag{7.6}
\end{equation}
如附录 C 所讨论，联络 $\omega$ 使我们能定义 $E$ 与 $\Phi^*VE$ 的截面的协变导数。与带规范场 4d A-模型一样，我们在局部用普通 1-形式表示这些协变导数：
\[
\begin{aligned}
\mathcal{D}\phi^I&=\dd\phi^I+\omega V^I\\
\mathcal{D}\chi^I&=\dd\chi^I+\Gamma^I{}_{JK}\dd\phi^J\chi^K+\omega\chi^K\nabla_KV^I
\end{aligned}
\]
定义在邻域 $U^{(\alpha,\gamma)}\subset\Sigma$ 上。

% [Tutor Note: p.90 式 (7.6) 可用 Gauss 曲率核对：$\omega=(\partial_2f)\dd\sigma^1-(\partial_1f)\dd\sigma^2$ 给出 $\dd\omega=-(\partial_1^2+\partial_2^2)f\,\dd\sigma^1\wedge\dd\sigma^2=-\Delta f\,\dd\sigma^1\wedge\dd\sigma^2$；而 $h_{\mu\nu}=\ee^{2f}\delta_{\mu\nu}$ 的 Gauss 曲率为 $K=-\ee^{-2f}\Delta f$，故 $\dd\omega=K\varepsilon$，即 $\dd\omega$ 恰为曲率 2-形式乘体积元。换成复坐标可写成 $\dd\omega=-2(\partial_z\partial_{\bar z}f)(\ii\dd z\wedge\dd\bar z)$，此式在 p.91 展开式中以 $+2(\partial_z\partial_{\bar z}f)H(\Phi)$ 的形式出现，两处符号可互相校验。]

令 $S_0$ 为把 A-模型作用量 (7.3) 中的导数简单替换为协变导数所得的最小耦合作用量。$S_0$ 的 BRST-正合部分仍 BRST 不变；然而在项 (7.4) 中把 $\dd\phi^I$ 换成 $\mathcal{D}\phi^I$ 会破坏其 BRST 不变性：
\begin{align*}
\delta_QS_0&=\delta_Q\Big(t\int_\Sigma\frac12\Omega_{IJ}\mathcal{D}\phi^I\wedge\mathcal{D}\phi^J\Big)\\
&=t\int_\Sigma\Big\{(\mathcal{L}_V\Omega)_{IJ}(\dd\phi^I\wedge\omega)\chi^J-\dd\big(\Omega_{IJ}\mathcal{D}\phi^I\chi^J\big)+\Omega_{IJ}V^I\chi^J\dd\omega\Big\}\\
&=t\int_\Sigma\Omega_{IJ}V^I\chi^J\dd\omega,
\end{align*}
其中我们用了 $\dd\Omega=0$。从第二行到第三行，我们用了 Kähler 形式 $\Omega$ 的 $U(1)$ 不变性，并丢掉了全导数项，因为按假定 $\Sigma$ 无边界。

为恢复 BRST 不变性，我们给作用量加一个额外（非最小）项 $S_1$，其 BRST 变分抵消上式。正如 §7.4 将说明的，项 $\Omega_{IJ}V^I\chi^J$ 的 BRST-正合性等价于 1-形式 $\iota_V\Omega$ 在 $X$ 上的 de Rham 上同调意义下恰当。事实上，我们认出这正是我们一开始所作的假定：群作用是 Hamilton 作用。整体矩映射 $H$ 的存在使我们可以写下表达式
\[
S_1=-t\int_\Sigma H(\Phi)\dd\omega,
\]
它具有恰当的 BRST 变分以抵消 $S_0$ 的变分。该项在 $\Sigma$ 上整体良定义，因为自旋联络的曲率 $\dd\omega$ 是整体定义的 2-形式，而函数 $H$ 在群作用的轨道上为常数；也就是说，
\[
H(\Phi^{(\alpha)})\dd\omega^{(\alpha)}=H(\Phi^{(\beta)})\dd\omega^{(\beta)}
\]

% [Tutor Note: p.90 三行推导的省略步：关键是 $\delta_Q$ 与协变导数可交换，即 $\delta_Q\mathcal{D}\phi^I=\mathcal{D}\delta_Q\phi^I=\mathcal{D}\chi^I$（$\omega$ 是由背景度规决定的自旋联络，与场无关，故 $\delta_Q\omega=0$）。在此基础上把 $\delta_Q(\tfrac12\Omega_{IJ}\mathcal{D}\phi^I\wedge\mathcal{D}\phi^J)$ 展开，用 $U(1)$ 不变性 $(\mathcal{L}_V\Omega)_{IJ}=0$ 消去一项，再把剩余部分写成全导数 $\dd(\Omega_{IJ}\mathcal{D}\phi^I\chi^J)$ 与 $\Omega_{IJ}V^I\chi^J\dd\omega$ 之和，即得第二行；全导数项在 $\partial\Sigma=\emptyset$ 时积分为零。这个恒等式也可用来独立核对 (7.7) 中 $\dd\omega$ 项的系数。]
% [Tutor Note: p.90 末 $S_1=-t\int_\Sigma H(\Phi)\dd\omega$ 的符号可按“抵消”反推核对：$\delta_QS_1=-t\int_\Sigma(\delta_QH)\dd\omega=-t\int_\Sigma\chi^I\partial_IH\,\dd\omega$，要与 $+t\int_\Sigma\Omega_{IJ}V^I\chi^J\dd\omega$ 相消，需 $\chi^I\partial_IH=+\Omega_{IJ}V^I\chi^J$。这与 (7.5) 印出的负号相矛盾（详见 p.88 批注），故本章的符号约定宜以“$S_1$ 必须抵消 $\delta_QS_0$”这一功能性要求为准。]

% [原文 p.91]
在重叠 $U^{(\alpha)}\cap U^{(\beta)}$ 上成立。总的来看，修正 A-模型的作用量由下式给出
\begin{equation}
S=\delta_QV+t\int_\Sigma\frac12\Omega_{IJ}\mathcal{D}\phi^I\wedge\mathcal{D}\phi^J-t\int_\Sigma H(\Phi)\dd\omega \tag{7.7}
\end{equation}
其中
\begin{equation}
V=-\frac{1}{4t}\int_\Sigma g_{i\bar j}\Big\{\rho_{\bar z}{}^i\big(\tilde P_z{}^{\bar j}+4\ii\mathcal{D}_z\phi^{\bar j}\big)+\tilde\rho_z{}^{\bar j}\big(P_{\bar z}{}^i-4\ii\mathcal{D}_{\bar z}\phi^i\big)\Big\}(\ii\dd z\wedge\dd\bar z). \tag{7.8}
\end{equation}
取 BRST 变分并在复坐标下写出，得到
\[
S=t\int_\Sigma\Big\{g_{i\bar j}\big(\mathcal{D}_z\phi^i\mathcal{D}_{\bar z}\phi^{\bar j}+\mathcal{D}_{\bar z}\phi^i\mathcal{D}_z\phi^{\bar j}+\ii\rho_{\bar z}{}^i\mathcal{D}_z\chi^{\bar j}-\ii\tilde\rho_z{}^{\bar j}\mathcal{D}_{\bar z}\chi^i\big)-\frac12R_{i\bar jk\bar l}\rho_{\bar z}{}^i\tilde\rho_z{}^{\bar j}\chi^k\chi^{\bar l}+2(\partial_z\partial_{\bar z}f)H(\Phi)\Big\}(\ii\dd z\wedge\dd\bar z).
\]
辅助场的运动方程为
\[
P_{\bar z}{}^i=2\ii\mathcal{D}_{\bar z}\phi^i,\qquad \tilde P_z{}^{\bar i}=-2\ii\mathcal{D}_z\phi^{\bar i}.
\]
修正 A-模型局域化在 A-模型全纯瞬子方程的协变化版本
\[
\mathcal{D}_{\bar z}\phi^i=\partial_{\bar z}\phi^i+\omega_{\bar z}V^i=0
\]
的解上。这个局部表达式等价于条件
\[
\dd\Phi\circ J_\Sigma=J_E\circ\dd\Phi,
\]
其中 $J_\Sigma$ 是 $\Sigma$ 上的复结构，$J_E$ 是由 $\omega$ 确定的 $E$ 的全空间上的复结构。（关于用底层主丛上的联络构造复结构 $J_E:TE\to TE$，见附录 C。）这个方程反过来又说明：映射 $\Phi:\Sigma\to E$ 是复流形之间的全纯映射。

最后注意，当 $\Sigma$ 是平坦环面时（此时自旋联络平凡），修正 A-模型的作用量退化为普通 A-模型的作用量。

% [Tutor Note: p.91 展开式中 $+2(\partial_z\partial_{\bar z}f)H(\Phi)$ 一项完全来自 $-t\int H\dd\omega$：由 (7.6) 得 $\dd\omega=-2(\partial_z\partial_{\bar z}f)(\ii\dd z\wedge\dd\bar z)$，故 $-t\int H\dd\omega=+2t\int(\partial_z\partial_{\bar z}f)H(\ii\dd z\wedge\dd\bar z)$。核对时另可用 $\partial_z\partial_{\bar z}=\tfrac14(\partial_1^2+\partial_2^2)$ 把此项与 $\dd\omega=-\Delta f\dd\sigma^1\wedge\dd\sigma^2$ 的结果对照，两条路径应给出同一系数，若差符号则说明 $(\ii\dd z\wedge\dd\bar z)$ 的定向约定被弄反。]
% [Tutor Note: p.91 的局域化方程 $\mathcal{D}_{\bar z}\phi^i=0$ 与 $\dd\Phi\circ J_\Sigma=J_E\circ\dd\Phi$ 的等价性依赖附录 C 中“用联络把水平/垂直分解搬到 $TE$”的构造：$\mathcal{D}_{\bar z}\phi^i=0$ 只保证 $\partial_{\bar z}$ 方向的像落在 $V^{1,0}E$ 里，是全纯截面的标准判据；当 $\omega$ 取平坦环面上的平凡联络（$f$ 为常数）时它退回 (7.4) 下方的 $\partial_{\bar z}\phi^i=0$，这正是正文最后一句“退化为普通 A-模型”的含义。]

\section{等变拓扑项}

考虑 BRST-非正合的部分
\[
S_{top}\equiv\int_\Sigma\frac12\Omega_{IJ}\mathcal{D}\phi^I\wedge\mathcal{D}\phi^J-\int_\Sigma H(\Phi)\dd\omega
\]
